\documentclass[12pt]{article}

% Layout.
\usepackage[top=1.2in, bottom=0.75in, left=1in, right=1in, headheight=1.0in, headsep=0pt]{geometry}

% Fonts.
\usepackage{mathptmx}
\usepackage[scaled=0.86]{helvet}
\renewcommand{\emph}[1]{\textsf{\textbf{#1}}}

% TiKZ.
\usepackage{tikz, pgfplots}
\usetikzlibrary{calc}
\pgfplotsset{my style/.append style={axis x line=middle, axis y line=middle, xlabel={$x$}, ylabel={$y$}}}
\pgfplotsset{compat=1.18}

% Misc packages.
\usepackage{amsmath,amssymb,latexsym}
\usepackage{graphicx}
\usepackage{array}
\usepackage{xcolor}
\usepackage{multicol}

% Commands to set various header/footer components.
\makeatletter
\def\doctitle#1{\gdef\@doctitle{#1}}
\doctitle{Use {\tt\textbackslash doctitle\{MY LABEL\}}.}
\def\docdate#1{\gdef\@docdate{#1}}
\docdate{Use {\tt\textbackslash docdate\{MY DATE\}}.}
\def\doccourse#1{\gdef\@doccourse{#1}}
\let\@doccourse\@empty
\def\docscoring#1{\gdef\@docscoring{#1}}
\let\@docscoring\@empty
\def\docversion#1{\gdef\@docversion{#1}}
\let\@docversion\@empty
\makeatother

% Headers and footers layout.
\makeatletter
\usepackage{fancyhdr}
\pagestyle{fancy}
\fancyhf{} % Clears all headers/footers.
\lhead{\emph{\@doctitle\hfill\@docdate} \medskip
\ifnum \value{page} > 1\relax\else\\
\emph{Name: \rule{3.5in}{1pt}\hfill \@docscoring}
\fi}

\rfoot{\emph{\@docversion}}
\lfoot{\emph{\@doccourse}}
\cfoot{\emph{\thepage}}
\renewcommand{\headrulewidth}{0pt}%
\makeatother

% Paragraph spacing
\parindent 0pt
\parskip 6pt plus 1pt

% A problem is a section-like command. Use \problem{5} for a problem worth 5 points.
\newcounter{probcount}
\newcounter{subprobcount}
\setcounter{probcount}{0}

\newcommand{\problem}[1]{%
\par
\addvspace{4pt}%
\setcounter{subprobcount}{0}%
\stepcounter{probcount}%
\makebox[0pt][r]{\emph{\arabic{probcount}.}\hskip1ex}\emph{[#1 points]}\hskip1ex}

\newcommand{\thesubproblem}{\emph{\alph{subprobcount}.}}

% like \problem but with name
\newcommand{\nameprob}[2]{%
\par
\addvspace{4pt}%
\setcounter{subprobcount}{0}%
\stepcounter{probcount}%
\makebox[0pt][r]{\emph{#1.}\hskip1ex}\emph{[#2 points]}\hskip1ex}

% Subproblems are an enumerate-like environment with a consistent
% numbering scheme.  Use \begin{subproblems}\item...\item...\end{subproblems}
\newenvironment{subproblems}{%
\begin{enumerate}%
\setcounter{enumi}{\value{subprobcount}}%
\renewcommand{\theenumi}{\emph{\alph{enumi}}}}%
{\setcounter{subprobcount}{\value{enumi}}\end{enumerate}}

% Blanks for answers in normal and math mode.
\newcommand{\blank}[1]{\rule{#1}{0.75pt}}
\newcommand{\mblank}[1]{\underline{\hspace{#1}}}
\def\emptybox(#1,#2){\framebox{\parbox[c][#2]{#1}{\rule{0pt}{0pt}}}}

% Misc.
\renewcommand{\d}{\displaystyle}
\newcommand{\ds}{\displaystyle}


\doctitle{Math 252: Quiz 4 }
\docdate{18 Sept 2025}
\doccourse{}
\docversion{}
\docscoring{\fbox{{\LARGE \strut}\blank{0.8in} / 25}}

\begin{document}
30 minutes maximum.  No aids (book, calculator, etc.) are permitted.  Show all work and use proper notation for full credit.  Answers should be in reasonably-simplified form.  25 points possible.

\begin{enumerate}
%[like 2.5 # 224,226]
\item (4 points) A 1-dimensional metal rod is 5 meters long (starting at $x=0$) and has density function $\rho(x)=e^{2x}$ kg/m.
	\begin{enumerate}
	\item (1 point) Compare the density of the rod at $x=1$ and at $x=4$. Where is the rod more dense? Explain your reasoning.
	\vfill
	\item (3 points) Find the mass of the 1-dimensional rod. Give units with your answer.
	\vfill
	\end{enumerate}
\newpage
% [Like 2.5 #236 exactly] 
\item (8 points) A spring has a natural length of $2$ meters. It requires 5 J of \textbf{work} to stretch the spring to a length of $2.5$ meters. 
	\begin{enumerate}
	\item (3 points) Find the spring constant $k$ in Hooke's Law.
	\vfill
	
	\vfill
	\item (2 points) Use your answer from part (a) to determine how much \textbf{force} the spring exerts if the spring is displaced exactly $0.5$ meters from its natural position. Include units with your answer.
	\vfill
	\item (3 points) How much \textbf{work} would it take to stretch the spring from $3$ meters to $4$ meters? Include units with your answer.
	\vfill
	\vfill
	\end{enumerate}
%like 2.6 # 251
\item (4 points) Calculate the center of mass for the collection of point masses given where $x_i$ gives the location of the point mass on the $x$-axis and $m_i$ is the mass. \\

\begin{tabular}{cc}
location&mass\\
\hline
$x_1=1$&$m_1=6$\\
$x_2=3$&$m_2=2$\\
$x_3=4$&$m_3=1$\\
\end{tabular}

\vspace{.3in}
\newpage
\item (9 points) Let $R$ be the region in the first quadrant bounded by $f(x)=1-x^2$. (Making a rough sketch of $R$ would probably be helpful here.)
	\begin{enumerate}
	\item (5 points) Find the $\overline{x},$ the $x$-coordinate of the \textbf{centroid} of $R.$ Recall that the \textbf{centroid} is the same as finding the \textbf{center of mass} when the density, $\rho,$ is 1. This will require evaluating two integrals, one to calculate $m$ and another for $M_y.$
	
	\vfill
	
	\vfill
	\item (2 points) \textbf{Set up but do not evaluate} the third integral you would need to determine $\overline{y}$, the $y$-coordinate of the centroid of $R.$ 
	\vfill
	\item (2 points) Suppose someone calculated $\overline{y} = 2/5.$ Is this value reasonable/plausible? Justify your conclusion.
	\vfill
	\end{enumerate}
\end{enumerate}
\end{document}

